DGFiP, REI 2017 et 2025, DVF 2023 à 2025 ; Insee, recensement 2022, comptes de patrimoine 2025 et Insee Analyses n° 91 ; calculs de l'auteur (synthese/calculs\_foncier.py, foncier/simulation.py).\par Note : proposition, toutes les composantes (taux national et nouvelle répartition communale) montant en ligne droite à partir de la première année, les droits étant divisés par deux dès l'annonce. Hausse annuelle : 90e et 99e centiles des logements classés selon leur hausse. Déficit : droits perdus non encore remplacés, sans reprise des ventes (avec une reprise de 6 \% par point : 18 \% de moins). Si les valeurs provisoires du sol n'existaient qu'au bout de trois ans, il faudrait ajouter 23 Md\euro{}.
